% Part: lambda-calculus
% Chapter: syntax
% Section: term-revisited

\documentclass[../../../include/open-logic-section]{subfiles}

\begin{document}

\olfileid{lam}{syn}{tr}
\olsection{$\alpha$-સામ્ય વર્ગો તરીકે પદો}

હવે પછી આપણે પદોને $\alpha$-સામ્ય સુધી લઈશું. એટલે કોઈ પદ
લખીએ ત્યારે તેનો $\alpha$-સામ્ય વર્ગ અભિપ્રેત હશે. દાખલા તરીકે,
$\lambd[a][\lambd[b][a c]]$ લખીએ ત્યારે આ પદ સાથે
$\alpha$-સામ્યવાળાં તમામ પદોનો ગણ સમજવાનો છે; તેમાં
$\lambd[a][\lambd[b][a c]]$ અને
$\lambd[b][\lambd[a][b c]]$ વગેરેનો સમાવેશ થાય છે.

અગાઉના વિભાગોમાં $N, Q$ વગેરે અક્ષરો પદ દર્શાવતા હતા;
હવે પછી તેઓ પદોના વર્ગો દર્શાવશે. આગળના અભ્યાસનો વિષય
પદોને બદલે આ વર્ગો હશે. $x, y$ વગેરે અક્ષરો હજી પણ
ચલો જ દર્શાવે છે.

વર્ગ~$M$ના કોઈપણ !!{element} માટે $\rep{M}$ સંકેત વાપરીશું.
એક કરતાં વધુ અવયવ જોઈતા હોય તો $\rep{M}[0], \rep{M}[1]$
વગેરે લખીશું.

વર્ણન સરળ બનાવવા પદો માટેના સંકેતો અહીં ફરી વાપરીશું.
વર્ગો પર નીચેની વ્યાખ્યા કરીએ:
\begin{defn}
  \begin{enumerate}
  \item $\lambd[x][N]$ એટલે $\lambd[x][\rep{N}]$ને
    સમાવતો વર્ગ.
  \item $PQ$ એટલે $\rep{P}\rep{Q}$ને સમાવતો વર્ગ.
  \end{enumerate}
\end{defn}

આ વ્યાખ્યાઓ સુવ્યાખ્યાયિત છે, કારણ કે $\alpha$-રૂપાંતરણ
રચના સાથે સુસંગત છે.

\begin{defn} \ollabel{def:fv}
  $\alpha$-સામ્ય વર્ગ~$M$ના \emph{મુક્ત ચલો}, એટલે
  $FV(M)$, એ $FV(\rep{M})$ તરીકે વ્યાખ્યાયિત થાય છે.
\end{defn}

આ વ્યાખ્યા સુવ્યાખ્યાયિત છે: \olref[alp]{thm:fv} મુજબ
$FV(\rep{M}[0]) = FV(\rep{M}[1])$.

વર્ગોમાં સ્થાનાપન્ન માટે પણ એ જ સંકેત વાપરીએ:
\begin{defn} \ollabel{def:sub}
  $M$માં $y$ની જગ્યાએ $R$નું \emph{સ્થાનાપન્ન}, એટલે
  $\Subst{M}{R}{y}$, એ એવા કોઈપણ $\rep{M}$ અને
  $\rep{R}$ માટે $\Subst{\rep{M}}{\rep{R}}{y}$ને
  સમાવતો $\alpha$-સામ્ય વર્ગ છે, જેના માટે અંદરનું
  સ્થાનાપન્ન વ્યાખ્યાયિત હોય.
  \sourcecorrection{OLLAM-027}{સ્થિર અંગ્રેજી મૂળ
  વર્ગો પરના સ્થાનાપન્નનું પરિણામ સીધું
  $\Subst{\rep{M}}{\rep{R}}{y}$ લખે છે, જે પદ છે,
  વર્ગ નથી. પરિણામ એ પદનો $\alpha$-સામ્ય વર્ગ છે
  એમ અહીં સ્પષ્ટ કર્યું છે.}
\end{defn}

આ પણ \olref[alp]{cor:sub}ના દાવા અનુસાર સુવ્યાખ્યાયિત
છે. જોકે એ ઉપસિદ્ધાંત માટેની અગાઉની સાબિતીમાં
\olref[alp]{thm:sub} અંગે નોંધાયેલો ગાળો હજી ખુલ્લો છે.

આ વ્યાખ્યાથી આપણી દલીલો ઘણી સરળ બને છે. દાખલા તરીકે:
\begin{align}
  \Subst{\lambd[x][x]}{y}{x} & \ollabel{eq:1}\\
  &= \Subst{\lambd[z][z]}{y}{x} \ollabel{eq:2}\\
  &= \lambd[z][\Subst{z}{y}{x}] \\
  &= \lambd[z][z]
\end{align}

આ સમીકરણો $\alpha$-સામ્ય વર્ગો વિશે છે: $\lambd[x][x]$
અને $\lambd[z][z]$ એક જ વર્ગમાં હોવાથી
\olref{eq:1}માંથી \olref{eq:2}માં જઈ શકાય છે. અગાઉ
સુધારેલી પદ-સ્તરની સ્થાનાપન્ન-વ્યાખ્યા મુજબ
$\Subst{\lambd[x][x]}{y}{x}$ વ્યાખ્યાયિત છે અને
તેનું મૂલ્ય $\lambd[x][x]$ છે; વર્ગ-સ્તરે તે
$\lambd[z][z]$ના વર્ગ બરાબર છે.
\sourcecorrection{OLLAM-028}{સ્થિર અંગ્રેજી મૂળ
\olref{eq:1}ના પદ-સ્તરના સ્થાનાપન્નને અવ્યાખ્યાયિત
કહે છે, પરંતુ અગાઉ સુધારેલા બદ્ધ-ચલ નિયમ પ્રમાણે તે
વ્યાખ્યાયિત અને યથાવત છે. અહીં પદની શાબ્દિક સમાનતા
અને વર્ગની સમાનતા વચ્ચેનો ભેદ બતાવ્યો છે. મૂળ
પ્રદર્શનની પહેલી પંક્તિમાં અંતે લટકતું વધારાનું
સમાનચિહ્ન પણ દૂર કર્યું છે.}

આ જ કારણથી આગળ સ્થાનાપન્ન માટે જરૂરી શરતો
સંતોષતા પ્રતિનિધિઓ પસંદ કર્યાં છે એમ માનીશું. દાખલા તરીકે,
$\Subst{\lambd[x][N]}{R}{y}$ દેખાય ત્યારે
$\lambd[x][N]$ એવો પ્રતિનિધિ પસંદ કરાયો છે કે
$x \neq y$ અને $x \notin FV(R)$.

$\lambd[x][x]$ને ``વર્ગ'' કહેવું થોડું વિચિત્ર લાગે છે.
અગાઉનાં $\lambda$-પદોથી ભેદ રાખવા આ વર્ગોને હવે પછી
$\Lambda$-પદો કહીશું; ભાગના બાકીના વિભાગોમાં ટૂંકમાં
``પદો'' પણ કહીશું.

\begin{editorial}
  પદોને આપણે હજી વિદાય આપી શકતા નથી: $\Lambda$-પદોની
  આખી વ્યાખ્યા $\lambda$-પદો પર આધારિત છે. વધુમાં,
  $\Lambda$-પદો પર વિધેયો વ્યાખ્યાયિત કરવાની સ્વતંત્ર
  રીત હજી આપી નથી. તેથી એવા વિધેયો પહેલાં
  $\lambda$-પદો પર વ્યાખ્યાયિત કરીને પછી, સ્થાનાપન્નની
  જેમ, $\Lambda$-પદો પર ``પ્રક્ષેપિત'' કરવા પડે.
  જોકે વાચક આવાં $\Lambda$-પદો પરના વિધેયો કેવી રીતે વ્યાખ્યાયિત થાય તે
  સાહજિક રીતે સમજી શકશે એમ માનીએ છીએ.
\end{editorial}
\end{document}
